\section{Empirical Design}
\label{sec:design}

\subsection{Predictive regression and inference}

All return tests use one timing convention: information at $t$ predicts the return at $t+1$. The general specification is
\begin{equation}
    Y_{t+1}=\alpha+bS_t+\Gamma'W_t+\varepsilon_{t+1},
    \label{eq:state_regression}
\end{equation}
where $W_t$ contains the relevant lagged benchmarks or controls. The headline conditional-mean difference sets $W_t$ empty. With a binary state, $b$ is the active-minus-inactive mean difference; monthly return coefficients are multiplied by twelve for reporting. In tables indexed by return month, equation~\eqref{eq:state_regression} is equivalently written with $S_{t-1}$.

The state's 92 active months form fifteen contiguous episodes, fourteen longer than one month. Treating all active months as independent would overstate the information in the sample. I report Newey--West statistics with 12 lags and, for binary-state tests, a wild episode bootstrap that assigns Rademacher weights to contiguous episodes and imposes the null. The episode bootstrap is the headline inference because continuation months do not create independent synchronized shocks. The estimates can be economically large while their independent common-state variation remains fifteen episodes.

Distributional diagnostics report medians, lower quantiles, skewness, expected shortfall, and conditional variance. A separate diagnostic removes each outcome's full-sample bottom-decile observations from both states before re-estimating equation~\eqref{eq:state_regression}. Because that deletion conditions on realized returns, it is not a preferred estimator. It answers only whether a handful of extreme observations mechanically generate the mean gap.

\subsection{Benchmarks and controls}

The benchmark hierarchy moves from nearby measurement alternatives to broader risk states. The raw current-inversion state asks whether fresh entry and confirmed release add information. The U.S. inversion asks whether the conventional single-country proxy spans the international state. The average G10 slope and its change ask whether the threshold is merely a nonlinear transformation of a smooth global slope. Nine country-specific live states ask whether synchronized crossings add information beyond their individual components.

Lagged NFCI and VIX controls address predictive spanning: do variables already known at $t$ absorb the state's forecast? Contemporaneous dollar and currency-volatility factors answer a different question: do those realized factors account for the loss after it occurs? Neither design identifies the shock represented by $S_t$. Joint specifications are reported without sequential orthogonalization because the ordering of closely related yield-curve transformations would otherwise determine which variable receives their shared variation.

\subsection{Exposure ordering}

The bilateral state coefficients come from
\begin{equation}
    \Delta s_{i,t+1}=a_i+b_iS_t+e_{i,t+1}.
    \label{eq:bilateral_state}
\end{equation}
The nine estimates $\widehat b_i$ are compared with currencies' average predetermined equity betas. The slope, fitted intercept at zero beta, and uncertainty describe the exposure gradient. A beta-sorted portfolio and a panel interaction $S_t\times\widehat\beta_{i,t}$ provide complementary tests: the portfolio is economically transparent; the panel uses monthly variation but still has only fifteen common-state episodes.

This design separates three claims. Prediction asks whether $S_t$ forecasts returns. Exposure accounting asks whether the realized losses are ordered by a predetermined market characteristic. Mechanism identification would require independent variation that distinguishes global growth, policy, term-premium, inflation, and intermediary channels. The first two are directly estimated; the third is not.

\subsection{Specification-family discipline}

Two different specification exercises appear in the paper and are never pooled. The first is the state-definition neighborhood reported for the baseline 10Y--2Y analysis. It varies breadth, release confirmation, tenor, and current-count repairs. These rows reveal where the historical result strengthens or attenuates, but they are selected displayed specifications rather than the complete set of rules examined. Their $p$-values are per specification; the joint resamples required for family adjustment are unavailable.

The second exercise is an independently executable public proxy. Its code declares and exhaustively reports 64 combinations of raw or live 10Y--3M states, breadth thresholds, release lengths, one- or two-month return lags, and target leg sizes of two or three, with all currencies tied at a cutoff included and equally weighted. It reports every coefficient, an inclusive circular-shift reference, and a common-calendar maximum-absolute-standardized-coefficient reference across the family. Under the maintained null that simultaneous cyclic shifts of the states relative to returns are exchangeable, this construction controls the declared family; without that assumption it is a finite rotation-reference diagnostic. The declaration was not preregistered before the data were examined. The public inputs and outcome differ from the baseline analysis, so this exercise measures sign prevalence and selection sensitivity rather than reproducing the headline estimate.

The evidentiary hierarchy follows from those designs. The baseline state supplies the economic result and cross-sectional ordering. Its displayed neighborhood shows known sensitivity. The public family shows what happens when a nearby declared rule family is exhaustively reported and jointly evaluated. A future observation under a dated, unchanged 10Y--2Y rule remains the clean prospective test.

\subsection{Claims and rejection patterns}

Four findings would run against the paper's interpretation: the distribution shift vanishing outside a few tail observations; the U.S. or average-slope measures spanning the configuration; no cross-sectional relation between losses and predetermined exposure; or a public nearby-rule family with no sign resemblance. The reported tests run against the first two patterns. The S\&P-based exposure tests run against the third, although the nine-point fit lacks a displayed confidence region, the panel interaction is imprecise, and the full-sample MSCI sort fails. The public evidence is weaker: negative coefficients are frequent, but no rule meets the 5-percent common-calendar rotation-reference criterion. That result preserves the cross-country-configuration question while ruling out a claim that the exact carry coefficient is invariant to reasonable measurement choices.
